\section*{CRS is not Metadata}

\begin{insight}
In many software systems, the coordinate reference system (CRS) is
treated as metadata.

It is stored, displayed, and occasionally converted — but rarely
considered part of the model itself.

\medskip

In railway alignment modelling, this view is fundamentally wrong.

\medskip

The CRS determines:
\begin{itemize}
\item how distances are measured,
\item how curvature is interpreted,
\item and how different datasets relate to each other.
\end{itemize}

\medskip

A shift of a few centimeters in coordinates may appear negligible.

A change in curvature is not.

\medskip

If coordinate systems are inconsistent:
\begin{itemize}
\item alignments no longer match,
\item curvature estimates become unstable,
\item and optimization results lose physical meaning.
\end{itemize}

\medskip

The problem is not numerical.

It is structural.

\medskip

A railway alignment does not exist independently of its spatial
reference.

The coordinate system is part of the model.
\end{insight}
